\chapter{2007年天津卷（文）}

\FigureLayoutDeclare{img/2007/tianjin_liberal/q16_fig1.png}{5.0cm}{7bp 10bp 7bp 0bp}

\section{单选题}

\begin{problem}
已知集合 \( S = \{x \in \R \mid x + 1 \geqslant 2\} \)，\( T = \{-2, -1, 0, 1, 2\} \)，则 \( S \cap T = \)（\quad）

\choices
    {\(\{2\}\)}
    {\(\{1, 2\}\)}
    {\(\{0, 1, 2\}\)}
    {\(\{-1, 0, 1, 2\}\)}
\end{problem}

\begin{answer}
B
\end{answer}

\begin{solution}
由 \(x+1\geqslant2\) 得 \(x\geqslant1\)，所以 \(S\cap T=\{1,2\}\)，故选 B.
\end{solution}

\begin{problem}
设变量 \(x\)，\(y\) 满足约束条件 \(\left\{\begin{aligned} x - y &\geqslant -1,\\ x + y &\leqslant 4,\\ y &\geqslant 2 \end{aligned}\right.\)，则目标函数 \(z = 2x + 4y\) 的最大值为（\quad）

\choices
    {\(10\)}
    {\(12\)}
    {\(13\)}
    {\(14\)}
\end{problem}

\begin{answer}
C
\end{answer}

\begin{solution}
约束条件等价于 \(y-1\leqslant x\leqslant4-y\)，且 \(y\geqslant2\)，故可行域的顶点为 \((1,2)\)，\((2,2)\)，\((\frac32,\frac52)\). 在这三个顶点处，目标函数 \(z=2x+4y\) 的值分别为 \(10\)，\(12\)，\(13\)，所以最大值为 \(13\)，故选 C.
\end{solution}

\begin{problem}
“ \(a = 2\) ”是“直线 \(ax + 2y = 0\) 平行于直线 \(x + y = 1\) ”的（\quad）

\choices
    {充分而不必要条件}
    {必要而不充分条件}
    {充分必要条件}
    {既不充分也不必要条件}
\end{problem}

\begin{answer}
C
\end{answer}

\begin{solution}
直线 \(ax+2y=0\) 的斜率为 \(-\frac a2\)，直线 \(x+y=1\) 的斜率为 \(-1\). 两直线平行当且仅当 \(-\frac a2=-1\)，即 \(a=2\)，因此“\(a=2\)”是题中结论的充分必要条件，故选 C.
\end{solution}

\begin{problem}
设 \(a = \log_{\frac12} 3\)，\(b = \left(\frac{1}{3}\right)^{0.2}\)，\(c = 2^{\frac{1}{3}}\)，则（\quad）

\choices
    {\(a < b < c\)}
    {\(c < b < a\)}
    {\(c < a < b\)}
    {\(b < a < c\)}
\end{problem}

\begin{answer}
A
\end{answer}

\begin{solution}
因为 \(a=\log_{\frac12}3<0\)，\(b=3^{-1/5}\) 满足 \(0<b<1\)，而 \(c=2^{1/3}>1\)，所以 \(a<b<c\)，故选 A.
\end{solution}

\begin{problem}
函数 \(y = \log_2(x + 4)\)（\(x > 0\)）的反函数是（\quad）

\choices
    {\(y = 2^{x} + 4\)（\(x > 2\)）}
    {\(y = 2^{x} + 4\)（\(x > 0\)）}
    {\(y = 2^{x} - 4\)（\(x > 2\)）}
    {\(y = 2^{x} - 4\)（\(x > 0\)）}
\end{problem}

\begin{answer}
C
\end{answer}

\begin{solution}
由 \(y=\log_2(x+4)\) 得 \(2^y=x+4\)，即 \(x=2^y-4\). 原函数的定义域为 \(x>0\)，所以其值域为 \(y>2\). 交换 \(x\)，\(y\) 得反函数 \(y=2^x-4\)，定义域为 \(x>2\)，故选 C.
\end{solution}

\begin{problem}
设 \(a\)，\(b\) 为两条直线，\(\alpha\)，\(\beta\) 为两个平面. 下列四个命题中，正确的命题是（\quad）

\choices
    {若 \(a\)，\(b\) 与 \(\alpha\) 所成的角相等，则 \(a \parallel b\)}
    {若 \(a \parallel \alpha\)，\(b \parallel \beta\)，\(\alpha \parallel \beta\)，则 \(a \parallel b\)}
    {若 \(a \subset \alpha\)，\(b \subset \beta\)，\(a \parallel b\)，则 \(\alpha \parallel \beta\)}
    {若 \(a \perp \alpha\)，\(b \perp \beta\)，\(\alpha \perp \beta\)，则 \(a \perp b\)}
\end{problem}

\begin{answer}
D
\end{answer}

\begin{solution}
两条直线与同一平面所成的角相等，不能推出两直线平行；分别平行于两个平行平面的两条直线也不一定平行；两个平面内的平行直线不能推出两个平面平行. 若 \(a\perp\alpha\)，\(b\perp\beta\)，且 \(\alpha\perp\beta\)，则 \(a\)，\(b\) 分别是两平面的法线，而互相垂直的两个平面的法线互相垂直，所以 \(a\perp b\)，故选 D.
\end{solution}

\begin{problem}
设双曲线 \(\frac{x^2}{a^2} - \frac{y^2}{b^2} = 1\)（\(a > 0\)，\(b > 0\)）的离心率为 \(\sqrt{3}\)，且它的一条准线与抛物线 \(y^2 = 4x\) 的准线重合，则此双曲线的方程为（\quad）

\choices
    {\(\frac{x^2}{12} - \frac{y^2}{24} = 1\)}
    {\(\frac{x^2}{48} - \frac{y^2}{96} = 1\)}
    {\(\frac{x^2}{3} - \frac{2y^2}{3} = 1\)}
    {\(\frac{x^2}{3} - \frac{y^2}{6} = 1\)}
\end{problem}

\begin{answer}
D
\end{answer}

\begin{solution}
设双曲线的半焦距为 \(c\)，则 \(c^2=a^2+b^2\). 由离心率 \(\frac ca=\sqrt3\)，得 \(b^2=c^2-a^2=2a^2\). 抛物线 \(y^2=4x\) 的准线为 \(x=-1\)，而双曲线的准线为 \(x=\pm\frac{a^2}{c}=\pm\frac a{\sqrt3}\)，所以 \(\frac a{\sqrt3}=1\)，即 \(a^2=3\)，\(b^2=6\). 因而双曲线方程为 \(\frac{x^2}{3}-\frac{y^2}{6}=1\)，故选 D.
\end{solution}

\begin{problem}
设等差数列 \(\{a_{n}\}\) 的公差 \(d\) 不为 \(0\)，\(a_{1} = 9d\). 若 \(a_{k}\) 是 \(a_{1}\) 与 \(a_{2k}\) 的等比中项，则 \(k =\)（\quad）

\choices
    {\(2\)}
    {\(4\)}
    {\(6\)}
    {\(8\)}
\end{problem}

\begin{answer}
B
\end{answer}

\begin{solution}
等差数列的通项为 \(a_n=a_1+(n-1)d=(n+8)d\). 由 \(a_k\) 是 \(a_1\) 与 \(a_{2k}\) 的等比中项，得 \(a_k^2=a_1a_{2k}\)，从而 \((k+8)^2d^2=9(2k+8)d^2\). 因为 \(d\ne0\)，所以 \(k^2-2k-8=0\)，即 \(k=4\) 或 \(k=-2\). 由于 \(k\in\N^*\)，故 \(k=4\)，选 B.
\end{solution}

\begin{problem}
设函数 \(f(x) = \left|\sin \left(x + \frac{\pi}{3}\right)\right|\)（\(x \in \R\)），则 \(f(x)\)（\quad）

\choices
    {在区间 \(\left[\frac{2\pi}{3}, \frac{7\pi}{6}\right]\) 上是增函数}
    {在区间 \(\left[-\pi, -\frac{\pi}{2}\right]\) 上是减函数}
    {在区间 \(\left[\frac{\pi}{8}, \frac{\pi}{4}\right]\) 上是增函数}
    {在区间 \(\left[\frac{\pi}{3}, \frac{5\pi}{6}\right]\) 上是减函数}
\end{problem}

\begin{answer}
A
\end{answer}

\begin{solution}
令 \(u=x+\frac\pi3\)，则 \(f(x)=|\sin u|\). 当 \(x\in\left[\frac{2\pi}{3},\frac{7\pi}{6}\right]\) 时，\(u\in\left[\pi,\frac{3\pi}{2}\right]\)，于是 \(f(x)=-\sin u\)，且 \(\cos u\leqslant0\)，所以 \(f'(x)=-\cos u\geqslant0\). 因此 \(f(x)\) 在该区间上是增函数，故选 A.
\end{solution}

\begin{problem}
设 \( f(x) \) 是定义在 \( \R \) 上的奇函数，且当 \( x \geqslant 0 \) 时，\( f(x) = x^2 \). 若对任意的 \( x \in [t, t + 2] \)，不等式 \( f(x + t) \geqslant 2f(x) \) 恒成立，则实数 \( t \) 的取值范围是（\quad）

\choices
    {\([\sqrt{2}, +\infty)\)}
    {\([2, +\infty)\)}
    {\(\left(0,2\right]\)}
    {\(\left[-\sqrt{2}, -1\right] \cup \left[\sqrt{2}, \sqrt{3}\right]\)}
\end{problem}

\begin{answer}
A
\end{answer}

\begin{solution}
由奇性，\(f(x)=\begin{cases}
-x^2, & x<0,\\
x^2, & x\geqslant0,
\end{cases}\) 因此 \(f\) 在 \(\R\) 上严格增，且对任意 \(x\) 有 \(2f(x)=f(\sqrt2x)\). 所以题中不等式等价于 \(x+t\geqslant\sqrt2x\)，即 \((\sqrt2-1)x\leqslant t\). 要使它对 \(x\in[t,t+2]\) 恒成立，只需且必须在最大值 \(x=t+2\) 处成立，于是

\[
(\sqrt2-1)(t+2)\leqslant t,
\]

解得 \(t\geqslant\sqrt2\)，故选 A.
\end{solution}

\section{填空题}

\begin{problem}
从一堆苹果中任取了 \(20\) 只，并得到它们的质量（单位：克）数据分布表如下：

\begin{center}
\begin{tabular}{|c|c|c|c|c|c|c|}
\hline
分组 & \([90,100)\) & \([100,110)\) & \([110,120)\) & \([120,130)\) & \([130,140)\) & \([140,150)\) \\
\hline
频数 & \(1\) & \(2\) & \(3\) & \(10\) & \(3\) & \(1\) \\
\hline
\end{tabular}
\end{center}

则这堆苹果中，质量不小于 \(120\text{ 克}\) 的苹果数约占苹果总数的\fillinblank{}\%.
\end{problem}

\begin{answer}
\(70\)
\end{answer}

\begin{solution}
质量不小于 \(120\text{ 克}\) 的苹果有 \(10+3+1=14\) 只，总数为 \(20\) 只，所以所占百分比为 \(\frac{14}{20}\times100\%=70\%\).
\end{solution}

\begin{problem}
\(\left(x + \frac{1}{x^2}\right)^9\) 的二项展开式中常数项是\fillinblank{}. （用数字作答）
\end{problem}

\begin{answer}
\(84\)
\end{answer}

\begin{solution}
二项展开式的通项为 \(\mathrm C_9^k x^{9-k}\left(\frac{1}{x^2}\right)^k=\mathrm C_9^k x^{9-3k}\). 令 \(9-3k=0\)，得 \(k=3\)，所以常数项为 \(\mathrm C_9^3=84\).
\end{solution}

\begin{problem}
一个长方体的各顶点均在同一球的球面上，且一个顶点上的三条棱的长分别为 \(1\)，\(2\)，\(3\)，则此球的表面积为\fillinblank{}.
\end{problem}

\begin{answer}
\(14\pi\).
\end{answer}

\begin{solution}
长方体的外接球直径等于长方体的体对角线，所以球的半径为 \(r=\frac12\sqrt{1^2+2^2+3^2}=\frac{\sqrt{14}}2\). 因而球的表面积为 \(4\pi r^2=4\pi\cdot\frac{14}{4}=14\pi\).
\end{solution}

\begin{problem}
已知两圆 \( x^{2} + y^{2} = 10 \) 和 \( (x - 1)^{2} + (y - 3)^{2} = 20 \) 相交于 \(A\)，\(B\) 两点，则直线 \(AB\) 的方程是\fillinblank{}.
\end{problem}

\begin{answer}
\(x+3y=0\)
\end{answer}

\begin{solution}
将第二个圆的方程展开，得 \(x^2+y^2-2x-6y-10=0\). 两圆的公共弦所在直线为两圆方程相减所得的根轴，因此由 \(x^2+y^2=10\) 得 \(2x+6y=0\)，即 \(AB\) 的方程为 \(x+3y=0\).
\end{solution}

\begin{problem}
在 \(\triangle ABC\) 中，\(AB = 2\)，\(AC = 3\)，\(D\) 是边 \(BC\) 的中点，则 \(\overrightarrow{AD} \cdot \overrightarrow{BC} =\fillinblank{}\).
\end{problem}

\begin{answer}
\(\frac52\)
\end{answer}

\begin{solution}
取 \(A\) 为向量起点，记 \(\overrightarrow{AB}=\bs b\)，\(\overrightarrow{AC}=\bs c\). 因为 \(D\) 是 \(BC\) 的中点，所以 \(\overrightarrow{AD}=\frac{\bs b+\bs c}{2}\)，且 \(\overrightarrow{BC}=\bs c-\bs b\). 因此
\[
\overrightarrow{AD}\cdot\overrightarrow{BC}=\frac{(\bs b+\bs c)\cdot(\bs c-\bs b)}2=\frac{|\bs c|^2-|\bs b|^2}{2}=\frac{AC^2-AB^2}{2}=\frac{9-4}{2}=\frac52.
\]
\end{solution}

\begin{problem}
如图，用 \(6\) 种不同的颜色给图中的 \(4\) 个格子涂色，每个格子涂一种颜色. 要求相邻的两个格子颜色不同，且两端的格子颜色也不同. 则不同的涂色方法共有\fillinblank{} 种. （用数字作答）

\bitmapfigure{img/2007/tianjin_liberal/q16_fig1.png}
\end{problem}

\begin{answer}
\(630\)
\end{answer}

\begin{solution}
先给最左边的格子选颜色，有 \(6\) 种；再给第二个格子选颜色，有 \(5\) 种. 对于第三个格子，若它与第一个格子同色，则有 \(5\) 种选法；若它与第一个格子不同色，则有 \(4\) 种选法，而第四个格子的颜色必须与第一个格子不同. 因而在固定前两个格子颜色后，后三个格子的有效选法数为 \(1\times5+4\times4=21\)，总数为 \(6\times5\times21=630\).
\end{solution}

\section{解答题}

\begin{problem}
在 \(\triangle ABC\) 中，已知 \(AC = 2\)，\(BC = 3\)，\(\cos A = -\frac{4}{5}\).

\begin{enumerate}
\item 求 \(\sin B\) 的值.
\item 求 \(\sin\left(2B+\frac{\pi}{6}\right)\) 的值.
\end{enumerate}
\end{problem}

\begin{solution}
\begin{enumerate}
\item 由 \(\cos A=-\frac45\) 且 \(A\) 为三角形内角，得 \(\sin A=\frac35\). 设 \(a=BC=3\)，\(b=CA=2\)，由正弦定理得 \(\frac{b}{\sin B}=\frac{a}{\sin A}\)，所以 \(\sin B=\frac{b\sin A}{a}=\frac25\). 又因为 \(A\) 为钝角，所以 \(B\) 为锐角，从而 \(\cos B=\sqrt{1-\sin^2B}=\frac{\sqrt{21}}5\).
\item \(\sin2B=2\sin B\cos B=\frac{4\sqrt{21}}{25}\)，\(\cos2B=1-2\sin^2B=\frac{17}{25}\). 因此 \(\sin\left(2B+\frac{\pi}{6}\right)=\sin2B\cos\frac{\pi}{6}+\cos2B\sin\frac{\pi}{6}=\frac{4\sqrt{21}}{25}\cdot\frac{\sqrt3}{2}+\frac{17}{25}\cdot\frac12=\frac{12\sqrt7+17}{50}\).
\end{enumerate}
\end{solution}

\begin{problem}
已知甲盒内有大小相同的 \(3\) 个红球和 \(4\) 个黑球，乙盒内有大小相同的 \(5\) 个红球和 \(4\) 个黑球. 现从甲，乙两个盒内各任取 \(2\) 个球.

\begin{enumerate}
\item 求取出的 \(4\) 个球均为红球的概率.
\item 求取出的 \(4\) 个球中恰有 \(1\) 个红球的概率.
\end{enumerate}
\end{problem}

\begin{solution}
\begin{enumerate}
\item 从甲盒取球的基本等可能结果数为 \(\mathrm C_7^2\)，从乙盒取球的基本等可能结果数为 \(\mathrm C_9^2\)，所以总数为 \(\mathrm C_7^2\mathrm C_9^2\). 取出的 \(4\) 个球均为红球的结果数为 \(\mathrm C_3^2\mathrm C_5^2\)，故所求概率为 \(P_1=\frac{\mathrm C_3^2\mathrm C_5^2}{\mathrm C_7^2\mathrm C_9^2}=\frac{30}{756}=\frac5{126}\).
\item 恰有 \(1\) 个红球包括甲盒取 \(1\) 个红球且乙盒取 \(0\) 个红球，或甲盒取 \(0\) 个红球且乙盒取 \(1\) 个红球. 因而所求概率为
\[
P_2=\frac{\mathrm C_3^1\mathrm C_4^1\mathrm C_4^2+\mathrm C_4^2\mathrm C_5^1\mathrm C_4^1}{\mathrm C_7^2\mathrm C_9^2}=\frac{72+120}{756}=\frac{16}{63}.
\]
\end{enumerate}
\end{solution}

\begin{problem}
如图，在四棱锥 \(P - ABCD\) 中，\(PA \perp\) 底面 \(ABCD\)，\(AB \perp AD\)，\(AC \perp CD\)，\(\angle ABC = 60^{\circ}\)，\(PA = AB = BC\)，\(E\) 是 \(PC\) 的中点.

\begin{enumerate}
\item 求 \(PB\) 和平面 \(PAD\) 所成的角的大小.
\item 证明 \(AE \perp\) 平面 \(PCD\).
\item 求二面角 \(A - PD - C\) 的大小.
\end{enumerate}

\bitmapfigure{img/2007/tianjin_liberal/q19_fig1.png}
\end{problem}

\begin{solution}
不妨设 \(PA=AB=BC=1\)，建立空间直角坐标系，使 \(A(0,0,0)\)，\(B(1,0,0)\)，\(P(0,0,1)\).
由 \(BC=1\) 及 \(\angle ABC=60^\circ\)，可取 \(C\left(\frac12,\frac{\sqrt3}{2},0\right)\). 由 \(AB\perp AD\)，设 \(D=(0,d,0)\). 又因为
\[
\overrightarrow{AC}\cdot\overrightarrow{CD}=\left(\frac12,\frac{\sqrt3}{2},0\right)\cdot\left(-\frac12,d-\frac{\sqrt3}{2},0\right)=0,
\]
所以 \(d=\frac2{\sqrt3}\)，即 \(D\left(0,\frac2{\sqrt3},0\right)\).

\begin{enumerate}
\item 平面 \(PAD\) 的方程为 \(x=0\)，而 \(\overrightarrow{PB}=(1,0,-1)\). 设 \(PB\) 与平面 \(PAD\) 所成的角为 \(\theta\)，则
\[
\sin\theta=\frac{|1|}{\sqrt{1^2+(-1)^2}}=\frac1{\sqrt2}.
\]
因此 \(\theta=\frac\pi4\).
\item \(E\) 是 \(PC\) 的中点，所以 \(E\left(\frac14,\frac{\sqrt3}{4},\frac12\right)\). 于是
\[
\overrightarrow{AE}\cdot\overrightarrow{PC}=\left(\frac14,\frac{\sqrt3}{4},\frac12\right)\cdot\left(\frac12,\frac{\sqrt3}{2},-1\right)=0,
\]
\[
\overrightarrow{AE}\cdot\overrightarrow{PD}=\left(\frac14,\frac{\sqrt3}{4},\frac12\right)\cdot\left(0,\frac2{\sqrt3},-1\right)=0.
\]
直线 \(PC\)，\(PD\) 是平面 \(PCD\) 内相交的两条直线，所以 \(AE\perp\) 平面 \(PCD\).
\item 取平面 \(PAD\) 的法向量 \(\bs n_1=(1,0,0)\)，平面 \(PCD\) 的法向量为
\[
\bs n_2=\overrightarrow{PC}\times\overrightarrow{PD}=\left(\frac1{2\sqrt3},\frac12,\frac1{\sqrt3}\right).
\]
设二面角 \(A-PD-C\) 的大小为 \(\varphi\)，则
\[
\cos\varphi=\frac{|\bs n_1\cdot\bs n_2|}{|\bs n_1||\bs n_2|}=\frac{\frac1{2\sqrt3}}{\sqrt{\frac1{12}+\frac14+\frac13}}=\frac{\sqrt2}{4}.
\]
所以二面角 \(A-PD-C\) 的大小为 \(\varphi=\arccos\frac{\sqrt2}{4}\).
\end{enumerate}
\end{solution}

\begin{problem}
在数列 \(\{a_{n}\}\) 中，\(a_{1} = 2\)，\(a_{n + 1} = 4a_{n} - 3n + 1\)，\(n \in \N^*\).

\begin{enumerate}
\item 证明数列 \(\{a_{n}-n\}\) 是等比数列.
\item 求数列 \(\{a_{n}\}\) 的前 \(n\) 项和 \(S_n\).
\item 证明不等式 \(S_{n+1} \leqslant 4S_n\) 对任意 \(n \in \N^*\) 皆成立.
\end{enumerate}
\end{problem}

\begin{solution}
\begin{enumerate}
\item 令 \(b_n=a_n-n\)，则 \(b_1=a_1-1=1\)，且
\[
b_{n+1}=a_{n+1}-(n+1)=4a_n-3n+1-n-1=4(a_n-n)=4b_n.
\]
所以数列 \(\{a_n-n\}\) 是首项为 \(1\)，公比为 \(4\) 的等比数列.
\item 由上知 \(a_n-n=b_n=4^{n-1}\)，所以 \(a_n=4^{n-1}+n\). 因而
\[
S_n=\sum_{k=1}^{n}\left(4^{k-1}+k\right)=\frac{4^n-1}{3}+\frac{n(n+1)}2.
\]
\item 由 \(S_{n+1}=S_n+a_{n+1}\)，不等式 \(S_{n+1}\leqslant4S_n\) 等价于 \(3S_n\geqslant a_{n+1}\). 又 \(a_{n+1}=4^n+n+1\)，故
\[
3S_n-a_{n+1}=4^n-1+\frac{3n(n+1)}2-(4^n+n+1)=\frac{(3n+4)(n-1)}2\geqslant0
\]
对任意 \(n\in\N^*\) 成立，所以 \(S_{n+1}\leqslant4S_n\).
\end{enumerate}
\end{solution}

\begin{problem}
设函数 \(f(x) = -x(x - a)^2\)（\(x \in \R\)），其中 \(a \in \R\).

\begin{enumerate}
\item 当 \(a = 1\) 时，求曲线 \(y = f(x)\) 在点 \((2, f(2))\) 处的切线方程.
\item 当 \(a \neq 0\) 时，求函数 \(f(x)\) 的极大值和极小值.
\item 当 \(a > 3\) 时，证明存在 \(k \in [-1,0]\)，使得不等式 \(f(k - \cos x) \geqslant f(k^2 - \cos^2 x)\) 对任意的 \(x \in \R\) 恒成立.
\end{enumerate}
\end{problem}

\begin{solution}
\begin{enumerate}
\item 当 \(a=1\) 时，\(f(x)=-x(x-1)^2\)，所以 \(f(2)=-2\)，且 \(f'(x)=-3x^2+4x-1\)，从而 \(f'(2)=-5\). 因此所求切线方程为 \(y+2=-5(x-2)\)，即 \(y=-5x+8\).
\item \(f'(x)=-(x-a)(3x-a)\)，所以驻点为 \(x=\frac a3\) 和 \(x=a\)，且
\[
f\left(\frac a3\right)=-\frac{4a^3}{27},\qquad f(a)=0.
\]
当 \(a>0\) 时，\(f\) 在 \(x=\frac a3\) 处取得极小值 \(-\frac{4a^3}{27}\)，在 \(x=a\) 处取得极大值 \(0\)；当 \(a<0\) 时，\(f\) 在 \(x=a\) 处取得极小值 \(0\)，在 \(x=\frac a3\) 处取得极大值 \(-\frac{4a^3}{27}\).
\item 取 \(k=-1\)，记 \(u=\cos x\)，则 \(u\in[-1,1]\)，并且
\[
k-u=-1-u\in[-2,0],\qquad k^2-u^2=1-u^2\in[0,1].
\]
当 \(a>3\) 且 \(z\leqslant1\) 时，\(z-a<0\)，\(3z-a<0\)，故 \(f'(z)=-(z-a)(3z-a)<0\)，即 \(f\) 在 \((-\infty,1]\) 上递减. 又
\[
(k^2-u^2)-(k-u)=1-u^2+1+u=(2-u)(u+1)\geqslant0.
\]
所以 \(k-u\leqslant k^2-u^2\)，从而 \(f(k-u)\geqslant f(k^2-u^2)\). 由于 \(k=-1\in[-1,0]\)，结论成立.
\end{enumerate}
\end{solution}

\begin{problem}
设椭圆 \(\frac{x^2}{a^2} + \frac{y^2}{b^2} = 1\)（\(a > b > 0\)）的左，右焦点分别为 \(F_1\)，\(F_2\)，\(A\) 是椭圆上的一点，\(AF_2 \perp F_1F_2\). 原点 \(O\) 到直线 \(AF_1\) 的距离为 \(\frac{1}{3}|OF_1|\).

\begin{enumerate}
\item 证明 \(a = \sqrt{2} b\).
\item 求 \(t \in (0, b)\) 使得下述命题成立：设圆 \(x^{2} + y^{2} = t^{2}\) 上任意点 \(M(x_{0}, y_{0})\) 处的切线交椭圆于 \(Q_{1}\)，\(Q_{2}\) 两点，则 \(OQ_{1} \perp OQ_{2}\).
\end{enumerate}
\end{problem}

\begin{solution}
\begin{enumerate}
\item 设椭圆的半焦距为 \(c\)，则 \(F_1=(-c,0)\)，\(F_2=(c,0)\)，且 \(c^2=a^2-b^2\). 由 \(AF_2\perp F_1F_2\)，设 \(A=(c,y_0)\). 将 \(A\) 代入椭圆方程，得
\[
\frac{c^2}{a^2}+\frac{y_0^2}{b^2}=1,
\]
从而 \(y_0^2=\frac{b^4}{a^2}\). 直线 \(AF_1\) 的方程为 \(y_0x-2cy+y_0c=0\)，所以原点到该直线的距离为 \(\frac{|cy_0|}{\sqrt{y_0^2+4c^2}}\). 由题意
\[
\frac{|cy_0|}{\sqrt{y_0^2+4c^2}}=\frac c3,
\]
得 \(c^2=2y_0^2=\frac{2b^4}{a^2}\). 结合 \(c^2=a^2-b^2\)，有
\[
a^4-a^2b^2-2b^4=0,
\]
即 \((a^2-2b^2)(a^2+b^2)=0\). 因为 \(a>b>0\)，所以 \(a^2=2b^2\)，即 \(a=\sqrt2b\).
\item 由（1）知椭圆方程化为 \(x^2+2y^2=2b^2\). 取圆 \(x^2+y^2=t^2\) 上任意一点 \(M(x_0,y_0)\)，则 \(x_0^2+y_0^2=t^2\)，且该点处的切线为 \(xx_0+yy_0=t^2\). 令 \(\bs d=(-y_0,x_0)\)，则切线上任意一点可表示为 \(M+s\bs d\). 设切线与椭圆的交点 \(Q_i=M+s_i\bs d\)，其中 \(s_1\)，\(s_2\) 是方程
\[
(x_0-sy_0)^2+2(y_0+sx_0)^2=2b^2
\]
的两个根. 由韦达定理，
\[
s_1s_2=\frac{x_0^2+2y_0^2-2b^2}{y_0^2+2x_0^2}.
\]
又因为 \(M\cdot\bs d=0\) 且 \(|\bs d|^2=t^2\)，所以
\[
\overrightarrow{OQ_1}\cdot\overrightarrow{OQ_2}=(M+s_1\bs d)\cdot(M+s_2\bs d)=t^2(1+s_1s_2).
\]
因此 \(OQ_1\perp OQ_2\) 当且仅当 \(s_1s_2=-1\)，这等价于
\[
\frac{x_0^2+2y_0^2-2b^2}{y_0^2+2x_0^2}=-1
\quad\Longleftrightarrow\quad
3(x_0^2+y_0^2)=2b^2
\quad\Longleftrightarrow\quad
3t^2=2b^2.
\]
所以所求 \(t=\sqrt{\frac23}\,b\). 该值属于 \((0,b)\)，且 \(t<b\) 时切线经过椭圆内部点，故对圆上的任意 \(M\) 都与椭圆交于两点，命题确实成立.
\end{enumerate}
\end{solution}
